% GENERATED FILE. Edit canonical lessons, then run npm run generate:books.
% canonical-source-hash: fnv1a64:443681fe13a53895
% canonical-lessons: IT-C325-principe, IT-C325-principessa, IT-C325-regina, IT-C325-regno, IT-C325-regnare

\chapter{Nel regno lontano}
\label{ch:it-nel-regno-lontano}
\hlchaptermodality{325}

\begin{chapteropening}
\textbf{By the end of this chapter:} \emph{I can tell a royal tale: the prince, the princess, the queen, the kingdom and how they reigned.}

Complete the last lesson of chapter 325: I can tell a royal tale: the prince, the princess, the queen, the kingdom and how they reigned.
\end{chapteropening}

\section[il principe]{il principe — the prince}

\label{lesson:IT-C325-principe}



\begin{quote}
Before the new one: say the Italian for the ending, then the Italian for the moral.
\end{quote}

\subsection*{You'll want to know: il principe}
\textbf{il principe} — \textquotedblleft{}the prince\textquotedblright{}.

\begin{quote}
Il principe è arrivato a cavallo.
\end{quote}

\begin{quote}
\emph{The prince arrived on horseback.}
\end{quote}

Plural \emph{i principi}.

\begin{cousinweb}
\emph{principe} is inherited from Latin \emph{prīnceps} \textquotedblleft{}the first, the chief\textquotedblright{} (EtymDB).
\end{cousinweb}

\begin{grammarlens}[title={What you've built}]
One more word for telling, explaining and describing at B1.
\end{grammarlens}

\subsection*{Your turn}
\begin{itemize}
\raggedright
  \item \emph{Say it:} \emph{il principe}
  \item \emph{Say it:} \emph{il principe}, once more
  \item \emph{Say it:} a sentence of your own with \emph{il principe}
  \item \emph{Recall:} say the Italian for the ending, then the Italian for the moral, then say \emph{il principe} again
\end{itemize}

\begin{tcolorbox}[breakable,colback=teal!4,colframe=teal!35!black,title={Before you move on}]
What does il principe mean? (\textquotedblleft{}The prince\textquotedblright{}.) Say it once more.
\end{tcolorbox}
\section[la principessa]{la principessa — the princess}

\label{lesson:IT-C325-principessa}



\begin{quote}
Before the new one: say the Italian for the moral, then the Italian for the prince.
\end{quote}

\subsection*{You'll want to know: la principessa}
\textbf{la principessa} — \textquotedblleft{}the princess\textquotedblright{}.

\begin{quote}
La principessa vive in un castello.
\end{quote}

\begin{quote}
\emph{The princess lives in a castle.}
\end{quote}

Built on \emph{principe} with \emph{-essa}.

\begin{cousinweb}
\emph{principessa} is \emph{principe} with \emph{-essa}, from Late Latin \emph{-issa}; the root is Latin \emph{prīnceps} \textquotedblleft{}chief\textquotedblright{} (EtymDB).
\end{cousinweb}

\begin{grammarlens}[title={What you've built}]
One more word for telling, explaining and describing at B1.
\end{grammarlens}

\subsection*{Your turn}
\begin{itemize}
\raggedright
  \item \emph{Say it:} \emph{la principessa}
  \item \emph{Say it:} \emph{la principessa}, once more
  \item \emph{Say it:} a sentence of your own with \emph{la principessa}
  \item \emph{Recall:} say the Italian for the moral, then the Italian for the prince, then say \emph{la principessa} again
\end{itemize}

\begin{tcolorbox}[breakable,colback=teal!4,colframe=teal!35!black,title={Before you move on}]
What does la principessa mean? (\textquotedblleft{}The princess\textquotedblright{}.) Say it once more.
\end{tcolorbox}
\section[la regina]{la regina — the queen}

\label{lesson:IT-C325-regina}



\begin{quote}
Before the new one: say the Italian for the prince, then the Italian for the princess.
\end{quote}

\subsection*{You'll want to know: la regina}
\textbf{la regina} — \textquotedblleft{}the queen\textquotedblright{}.

\begin{quote}
La regina ha tre figli.
\end{quote}

\begin{quote}
\emph{The queen has three children.}
\end{quote}

Also the queen in chess.

\begin{cousinweb}
\emph{regina} is inherited from Latin \emph{rēgīna} \textquotedblleft{}queen\textquotedblright{} (EtymDB), built on \emph{rēx} \textquotedblleft{}king\textquotedblright{}.
\end{cousinweb}

\begin{grammarlens}[title={What you've built}]
One more word for telling, explaining and describing at B1.
\end{grammarlens}

\subsection*{Your turn}
\begin{itemize}
\raggedright
  \item \emph{Say it:} \emph{la regina}
  \item \emph{Say it:} \emph{la regina}, once more
  \item \emph{Say it:} a sentence of your own with \emph{la regina}
  \item \emph{Recall:} say the Italian for the prince, then the Italian for the princess, then say \emph{la regina} again
\end{itemize}

\begin{tcolorbox}[breakable,colback=teal!4,colframe=teal!35!black,title={Before you move on}]
What does la regina mean? (\textquotedblleft{}The queen\textquotedblright{}.) Say it once more.
\end{tcolorbox}
\section[il regno]{il regno — the kingdom; the reign}

\label{lesson:IT-C325-regno}



\begin{quote}
Before the new one: say the Italian for the princess, then the Italian for the queen.
\end{quote}

\subsection*{You'll want to know: il regno}
\textbf{il regno} — \textquotedblleft{}the kingdom; the reign\textquotedblright{}.

\begin{quote}
Il regno era piccolo ma ricco.
\end{quote}

\begin{quote}
\emph{The kingdom was small but rich.}
\end{quote}

Also the time a ruler reigns: \emph{durante il suo regno}.

\begin{cousinweb}
\emph{regno} is inherited from Latin \emph{rēgnum} \textquotedblleft{}royal power, a kingdom\textquotedblright{} (EtymDB).
\end{cousinweb}

\begin{grammarlens}[title={What you've built}]
One more word for telling, explaining and describing at B1.
\end{grammarlens}

\subsection*{Your turn}
\begin{itemize}
\raggedright
  \item \emph{Say it:} \emph{il regno}
  \item \emph{Say it:} \emph{il regno}, once more
  \item \emph{Say it:} a sentence of your own with \emph{il regno}
  \item \emph{Recall:} say the Italian for the princess, then the Italian for the queen, then say \emph{il regno} again
\end{itemize}

\begin{tcolorbox}[breakable,colback=teal!4,colframe=teal!35!black,title={Before you move on}]
What does il regno mean? (\textquotedblleft{}The kingdom; the reign\textquotedblright{}.) Say it once more.
\end{tcolorbox}
\section[regnare]{regnare — to reign}

\label{lesson:IT-C325-regnare}



\begin{quote}
Before the new one: say the Italian for the queen, then the Italian for the kingdom.
\end{quote}

\subsection*{You'll want to know: regnare}
\textbf{regnare} — \textquotedblleft{}to reign\textquotedblright{}.

\begin{quote}
Ha regnato per quarant'anni.
\end{quote}

\begin{quote}
\emph{She reigned for forty years.}
\end{quote}

Also figurative: \emph{regna il silenzio}.

\begin{cousinweb}
\emph{regnare} is inherited from Latin \emph{rēgnāre} \textquotedblleft{}to reign\textquotedblright{} (EtymDB).
\end{cousinweb}

\begin{grammarlens}[title={What you've built}]
One more word for telling, explaining and describing at B1.
\end{grammarlens}

\subsection*{Your turn}
\begin{itemize}
\raggedright
  \item \emph{Say it:} \emph{regnare}
  \item \emph{Say it:} \emph{regnare}, once more
  \item \emph{Say it:} a sentence of your own with \emph{regnare}
  \item \emph{Recall:} say the Italian for the queen, then the Italian for the kingdom, then say \emph{regnare} again
\end{itemize}

\begin{tcolorbox}[breakable,colback=teal!4,colframe=teal!35!black,title={Before you move on}]
What does regnare mean? (\textquotedblleft{}To reign\textquotedblright{}.) Say it once more.
\end{tcolorbox}
